% ---------------------------------------------------------------------------
% Preamble -- Time Series Term Project, Tunis Business School
% ---------------------------------------------------------------------------
\documentclass[11pt,a4paper]{article}

% ---------------- encoding, language, fonts -------------------------------
\usepackage[T1]{fontenc}
\usepackage[utf8]{inputenc}
\usepackage[english]{babel}
\usepackage{lmodern}
\usepackage{microtype}

% ---------------- page geometry -------------------------------------------
\usepackage[a4paper,top=2.5cm,bottom=2.5cm,left=2.5cm,right=2.5cm]{geometry}
\setlength{\parskip}{0.35em}
\setlength{\parindent}{0pt}

% ---------------- mathematics ---------------------------------------------
\usepackage{amsmath,amssymb,amsthm,mathtools}
\usepackage{bm}

% ---------------- tables, figures, lists ----------------------------------
\usepackage{booktabs}
\usepackage{array}
\usepackage{multirow}
\usepackage{longtable}
\usepackage{graphicx}
\usepackage{float}
\usepackage[font=small,labelfont=bf,justification=raggedright,singlelinecheck=false]{caption}
\usepackage{enumitem}
\usepackage{rotating}

% ---------------- typography and colour -----------------------------------
\usepackage[table]{xcolor}
\usepackage{titlesec}
\usepackage{fancyhdr}
\usepackage[colorlinks=true,linkcolor=black,citecolor=black,urlcolor=blue!55!black]{hyperref}
\usepackage{natbib}
\bibliographystyle{apalike}

% ---------------- floats that stay put in the pandoc fallback ------------
\graphicspath{{../reports/figures/}}
\newcommand{\figdir}{../reports/figures}
\newcommand{\tabdir}{../reports/tables}

% ---------------- headers / footers ---------------------------------------
\pagestyle{fancy}
\fancyhf{}
\fancyhead[L]{\footnotesize Time Series Analysis and Forecasting --- Term Project}
\fancyhead[R]{\footnotesize\thepage}
\renewcommand{\headrulewidth}{0.4pt}

\titleformat{\section}{\large\bfseries}{\thesection}{0.7em}{}
\titleformat{\subsection}{\normalsize\bfseries}{\thesubsection}{0.6em}{}

% ---------------- notation shortcuts --------------------------------------
\newcommand{\E}{\mathbb{E}}
\newcommand{\Var}{\mathrm{Var}}
\newcommand{\Cov}{\mathrm{Cov}}
\newcommand{\eps}{\varepsilon}
\newcommand{\B}{B}
\newcommand{\yhat}{\hat{y}}
\newcommand{\pihi}{\phi}
\DeclareMathOperator*{\argmin}{arg\,min}
\DeclareMathOperator*{\argmax}{arg\,max}

% ---------------- report metadata ----------------------------------------
\newcommand{\reporttitle}{Time Series Modeling and Forecasting of US Consumer Price Inflation}
\newcommand{\reportsubtitle}{A complete Box--Jenkins analysis of the Kaggle US CPI-U dataset, 1957--2020}
\newcommand{\course}{Time Series Analysis and Forecasting}
\newcommand{\institution}{Tunis Business School --- Finance / Business Administration}
\newcommand{\academicyear}{Academic year 2026/2027}


\begin{document}

% ===========================================================================
%  TITLE PAGE
% ===========================================================================
\begin{titlepage}
\centering
\vspace*{1.2cm}
{\scshape\large \institution \par}
\vspace{0.4cm}
{\scshape \course \par}
\vspace{2.0cm}
\rule{\textwidth}{0.6pt}\\[0.6cm]
{\LARGE\bfseries \reporttitle \par}
\vspace{0.5cm}
{\large \reportsubtitle \par}
\\[0.6cm]\rule{\textwidth}{0.6pt}

\vspace{1.8cm}
\begin{tabular}{r@{\hspace{1.2em}}l}
\textbf{Forecasting question} & One- to five-year ahead US inflation \\
\textbf{Data} & Kaggle ``US Inflation Data'' (CPI-U, monthly, 1913--2020) \\
\textbf{Estimation sample} & January 1957 -- December 2020 ($n=768$) \\
\textbf{Method} & Box--Jenkins: SARIMA identification, estimation, diagnostics \\
\textbf{Software} & Python 3.11 (\texttt{pandas}, \texttt{statsmodels}, \texttt{matplotlib}) \\
\end{tabular}

\vfill
{\large \academicyear \par}
\vspace{0.3cm}
{\normalsize Submitted as the term project of the course \course \par}
\end{titlepage}

% ===========================================================================
\pagenumbering{roman}
\section*{Executive summary}
\addcontentsline{toc}{section}{Executive summary}

This project asks a single decision-relevant question: \emph{given the history of US
consumer prices, what inflation rate should a Tunisian bank, treasurer or importer
budget for over the next one to five years, and how much confidence can be placed
in that number?} The answer is built from the full Box--Jenkins methodology on the
monthly US Consumer Price Index for All Urban Consumers (CPI-U, seasonally
adjusted), taken from a public Kaggle dataset covering January 1913 to December
2020 and estimated on the post-1957 sample ($n=768$).

The empirical sequence is the one taught in the course. The price index is
non-stationary: its autocorrelation function decays almost linearly and its
variance grows with its level, and the augmented Dickey--Fuller and KPSS tests
agree that one regular difference is required ($d=1$). The Box--Cox profile
likelihood points to a power transform between the log and the square root; the
logarithm is retained because it converts the index into a growth rate and because
every conclusion is unchanged under either transform. The differenced series,
monthly inflation $\pi_t = 100\Delta\ln P_t$, is stationary, but its autocorrelation
function shows a clear annual echo at lags 11--13, which justifies a
\emph{multiplicative seasonal} model rather than a purely non-seasonal one.

Model selection over a grid of multiplicative SARIMA specifications
($p,q\in\{0,\dots,3\}$, $(P,Q)\in\{(0,0),(1,0),(0,1),(1,1)\}$, $d=1$, $D=0$)
retains, on the joint criterion of lowest BIC among specifications whose residuals
pass the portmanteau tests,

\begin{equation}
(1-\phi_1 B)(1-\Phi_1 B^{12})\,\Delta y_t
\;=\; c \;+\; (1+\theta_1 B+\theta_2 B^2+\theta_3 B^3)(1+\Theta_1 B^{12})\,\eps_t ,
\label{eq:final}
\end{equation}
i.e.\ ARIMA$(1,1,3)(1,0,1)_{12}$ with $y_t=100\ln P_t$. All eight parameters except the constant are significant at
the 1\% level, the model is stationary and invertible, and $\hat\sigma_\eps=0.259$
percentage points of monthly log inflation. Of the 63 specifications estimated,
54 fail the Ljung--Box test at 12 lags; the selected one does not (12 lags:
$p=0.21$; 24 lags: $p=0.10$).

Out of sample, the model is re-estimated on 1957--2015 and used to forecast the
60 months 2016--2020. It attains an RMSE of $1.79$ on the log index against
$6.23$ for a random walk, $4.49$ for a random walk with drift and $1.86$ for a
seasonal random walk, a $71\%$ reduction in RMSE relative to the naive benchmark.
On the index itself the MAPE is $1.59\%$ over five years. The 95\% interval
covers the realised path throughout, and the one-step-ahead rolling evaluation
(RMSE $0.225$) is close to the theoretical innovation standard deviation
($\sigma=0.261$), which is the signature of a correctly specified conditional mean.

In decision terms: the model puts long-run US inflation at about
$2.9\%$ a year, implies a half-life of inflation shocks of roughly 16 months, and
would have predicted 2020 inflation at $1.7\%$ against a realisation of $1.35\%$.
The two episodes it cannot capture are the oil-embargo spike of 1973 and the
deflationary collapse of late 2008; both show up as large residuals and both are
the source of the volatility clustering that a pure ARIMA model leaves unexplained.

\vspace{0.6em}
\noindent\textbf{Keywords:} Box--Jenkins methodology; SARIMA; CPI-U; inflation forecasting;
stationarity; Ljung--Box test; forecast evaluation.

\clearpage
\tableofcontents
\clearpage
\pagenumbering{arabic}

% ===========================================================================
\section{Introduction}
% ===========================================================================

\subsection{The forecasting question}

A price index is one of the few economic series whose future value every
business has to budget for. A Tunisian importer quoting in US dollars, a bank
pricing a dollar-denominated loan, or a treasury desk setting a nominal return
all need a number for US inflation over the life of the contract, together with
an honest statement of the uncertainty around it. The forecasting question of
this project is therefore:

\begin{quote}
\emph{What rate of US consumer price inflation should be expected over horizons of
one month to five years, given the history of the CPI up to December 2020, and how
wide is the interval within which the realised rate should fall?}
\end{quote}

The horizon matters. A one-month-ahead forecast is useful for indexation clauses;
a five-year-ahead forecast is what a corporate bond or a long-term supply contract
requires. The forecast interval is as important as the point forecast, because a
financing decision depends on the whole distribution, not on its mean.

\subsection{Why a univariate time-series model}

The course covers univariate Box--Jenkins modelling, and the question is well
matched to it for three reasons. First, inflation persistence is a genuine
time-series phenomenon: this month's inflation tells us a great deal about next
month's, and a correctly specified ARMA structure exploits that dependence.
Second, parsimony has practical value: a model with eight parameters can be
re-estimated monthly by a junior analyst, which a large macro-econometric system
cannot. Third, and most importantly for a course in forecasting, the framework
produces a \emph{quantified} forecast error variance, so the interval around the
forecast is derived rather than asserted.

The obvious limitation, stated at the outset, is that a univariate model has
nothing to say about the \emph{causes} of inflation --- money growth, import
prices, output gaps. It extrapolates the process that generated the past, and it
will fail at genuine regime changes. Section~\ref{sec:limits} returns to this
point.

\subsection{Structure of the report}

Section~\ref{sec:data} documents the dataset, the sampling decisions and the
descriptive statistics. Section~\ref{sec:components} treats the components of the
series and the transformations of Chapter~2, including the Box--Cox analysis, the
ACF/PACF identification and the unit-root tests that fix $d$. Section~\ref{sec:id}
presents the identification and estimation of Chapter~6 and the three estimation
methods of Chapter~4. Section~\ref{sec:diag} reports the diagnostic checks.
Section~\ref{sec:fcst} develops the MMSE forecast recursion of Chapter~5 and
evaluates the forecasts against a hold-out sample and three benchmarks.
Section~\ref{sec:interpretation} translates the statistics into business language,
Section~\ref{sec:limits} lists the limitations, and Section~\ref{sec:conclusion}
concludes. The code that produces every number and figure in this report is in the
repository described in Appendix~\ref{app:code}.

% ===========================================================================
\section{Data: source, frequency, sample and descriptive statistics}
\label{sec:data}

\subsection{Source and frequency}

The data are the Kaggle dataset
\emph{``US Inflation Data (updated till May 2021)''}%
\footnote{\url{https://www.kaggle.com/datasets/varpit94/us-inflation-data-updated-till-may-2021}.},
distributed as a single CSV file with two columns, \texttt{Yearmon} (a date in
\texttt{dd-mm-yyyy} form, always the first of the month) and \texttt{CPI} (the
index value). It contains the monthly US Consumer Price Index for All Urban
Consumers, series CPIAUCSL on the Federal Reserve's FRED database, i.e.\ the
seasonally adjusted index for all items with base 1982--84 $=100$.

The file holds 1{,}296 monthly observations from 1913-01 to 2020-12, with no
missing month, no duplicated month and no missing value; the acquisition script
verifies the SHA-256 digest of the file against the archived copy
(\texttt{d54f79b8\dots3ece33c}) so the analysis is reproducible from the raw CSV.
Table~\ref{tab:dataaudit} summarises the audit.

\begin{table}[htbp]
\centering
\small
\caption{Data audit of the raw Kaggle file and of the estimation sample.}
\label{tab:dataaudit}
\begin{tabular}{lrr}
\toprule
Item & Raw file & Estimation sample \\
\midrule
Frequency & monthly & monthly \\
First observation & 1913-01 & 1957-01 \\
Last observation & 2020-12 & 2020-12 \\
Observations & 1{,}296 & 768 \\
Missing months & 0 & 0 \\
Duplicate months & 0 & 0 \\
Minimum day gap between records & 28 & 28 \\
Maximum day gap between records & 31 & 31 \\
\bottomrule
\end{tabular}
\end{table}

\subsection{Why the sample starts in 1957}

The dataset begins in 1913, but the estimation sample starts in 1957-01 for two
documented reasons. The pre-1957 CPI-U series is a \emph{backward extrapolation}:
the modern index was only linked to the earlier cost-of-living series in 1953, so
the pre-1957 values are constructed rather than collected, and their dynamics
reflect the construction method as much as the economy. In addition, the
immediate post-war observations (1946--1948 and 1950--1951) contain one-off
releases of wartime price controls, movements that are not draws from the same
stochastic process as peacetime inflation.

The cost of this choice is that the project says nothing about the deflation of
1921 or 1930--1933. That is acceptable: the forecasting question concerns the
inflation process that a contract written today will face, and modern US inflation
has never returned to the interwar regime. Figure~\ref{fig:history} shows the full
sample with the estimation window shaded, so the reader can see exactly what is
excluded.

\begin{figure}[htbp]
\centering
\includegraphics[width=\textwidth]{\figdir/f01_series_components.png}
\caption{The three series used in the analysis. (a) The CPI level is dominated by a
near-deterministic upward trend and its variance grows with the level, the classic
signature of a series that must be transformed and differenced. (b) Taking the
logarithm stabilises the variance and linearises the trend. (c) Monthly and annual
inflation fluctuate around a positive mean with long swings (the cyclical
component) and no trend.}
\label{fig:history}
\end{figure}

\subsection{Descriptive statistics}

Table~\ref{tab:descriptives} reports the moments of the four working series. Note
the contrast that drives everything that follows: the level and the log level have
a standard deviation nearly equal to their own mean (76.6 against 125.7, and 76.3
against 458.5), because both are trending, while monthly inflation has a mean of
$0.293\%$ and a standard deviation of $0.352\%$ --- a coefficient of variation
above unity but no trend at all. Monthly inflation is very mildly right-skewed
(0.03) and strongly leptokurtic (excess kurtosis 3.08), so large monthly moves are
far more common than a Gaussian model would suggest. Annual inflation, by
contrast, is strongly right-skewed (1.56), because the 1970s produced double-digit
rates while the floor is bounded below near zero.

\input{\tabdir/t01_descriptives}

Figure~\ref{fig:decades} makes the regime structure explicit. Annual inflation
averaged $2.3\%$ in the 1960s, $6.8\%$ in the 1970s and $5.4\%$ in the 1980s, fell
to $2.95\%$ in the 1990s and $2.5\%$ in the 2000s, and reached $1.75\%$ in the
2010s (Figure~\ref{fig:decades}). This is not a stationary mean: it is a sequence of regimes, and the largest
of them --- the Volcker disinflation, which took annual inflation from a peak of
$13.8\%$ in March 1980 to $3.8\%$ by the end of 1982, and below $3\%$ a year later
--- is the single most influential event in the sample. Section~\ref{sec:diag} shows that the residuals of the final model
are dominated by exactly this transition.

\begin{figure}[htbp]
\centering
\includegraphics[width=\textwidth]{\figdir/f17_decades.png}
\caption{Distribution of annual inflation by decade. The boxes are the interquartile
ranges, the diamonds the means; the dispersion, not only the level, changes across
decades, and the 1970s and 1980s dominate the sample mean.}
\label{fig:decades}
\end{figure}

\begin{table}[htbp]
\centering
\small
\caption{Annual CPI inflation by decade (twelve-month log change, per cent). The
1950s and 2020s columns are partial decades (24 and 12 months).}
\label{tab:decades}
\begin{tabular}{lrrrr}
\toprule
Decade & Mean & S.d. & Minimum & Maximum \\
\midrule
1950s (24 obs.) & 1.85 & 1.05 & 0.35 & 3.56 \\
1960s & 2.30 & 1.44 & 0.67 & 6.01 \\
1970s & 6.82 & 2.53 & 2.67 & 12.48 \\
1980s & 5.36 & 3.27 & 1.09 & 13.76 \\
1990s & 2.95 & 1.08 & 1.37 & 6.10 \\
2000s & 2.53 & 1.41 & $-2.12$ & 5.45 \\
2010s & 1.75 & 0.85 & $-0.20$ & 3.80 \\
2020s (12 obs.) & 1.23 & 0.69 & 0.12 & 2.46 \\
\bottomrule
\end{tabular}
\end{table}

% ===========================================================================
\section{Components, transformations and stationarity (Chapters 1--2)}
\label{sec:components}

\subsection{Reading the components off the plot}

Chapter~1 asks for the trend, seasonal, cyclical and irregular components to be
identified before any formal modelling.

\textbf{Trend.} The log index rises almost monotonically from 331.8 in 1957-01 to
556.3 in 2020-12, a total log change of 224.5 points. Spread over 64 years that
is $3.51$ points per year --- an average annual inflation rate of $3.5\%$ ---
while the least-squares slope of the log index is $4.03$ points per year, so the
average was dragged down by the slow decades at the start and the end of the
sample. The trend is not
linear in the level but is close to linear in the log, which is itself evidence
for the logarithmic transformation.

\textbf{Cyclical component.} The deviations from trend are long and slow. Annual
inflation rises for roughly a decade (1963--1974 and 1976--1980), stays high for
several years, then falls for a decade (1981--1992, 2008--2015). A cycle of this
length is a business-cycle phenomenon and cannot be captured by seasonal terms; it
becomes part of the ARMA dynamics.

\textbf{Seasonal component.} This is the interesting case, because the CPI-U
series in the Kaggle file is \emph{seasonally adjusted} by the source. A classical
STL decomposition of $\pi_t$ (Figure~\ref{fig:stl}) attributes $37.6\%$ of the
variance of monthly inflation to trend, $24.4\%$ to a seasonal component with a
peak-to-trough amplitude of $1.11$ percentage points, and $41.4\%$ to the
remainder. The seasonal share is therefore not zero even though the series is
adjusted. The explanation is that the adjustment removes \emph{deterministic}
calendar effects but cannot remove the fact that price resets cluster at the turn
of the year; what remains is a stochastic annual echo, and it is visible in the
autocorrelation function at lags 11--13 (Section~\ref{sec:acf}). This is precisely
the kind of residual seasonality that the seasonal SARMA models of Chapter~7 are
designed to absorb, and it is why the final model carries a
$\Theta(B^{12})$ term.

\begin{figure}[htbp]
\centering
\includegraphics[width=\textwidth]{\figdir/f02_stl_decomposition.png}
\caption{Additive STL decomposition of monthly inflation. The trend line shows the
Great Inflation and the Great Moderation; the seasonal line is small but regular;
the remainder contains the oil-embargo and financial-crisis shocks.}
\label{fig:stl}
\end{figure}

\textbf{Irregular component.} Eight months have a standardised monthly inflation
above three in absolute value: 1973-08 (the oil embargo, $+1.79\%$), January to
March 1980 (the second oil shock and the imposition of credit controls),
2005-11, and October to December 2008 (the financial crisis, with 2008-11 at
$-1.93\%$, the largest single monthly decline in the sample).

\subsection{Variance-stabilising transformation and the Box--Cox family}

Chapter~2 requires a variance-stabilising transformation. The power transform
$g(P_t;\lambda)$ of the CPI level and its profile log-likelihood were evaluated
over $\lambda\in[-0.5,1.5]$, using
\begin{equation}
g(P;\lambda) = \begin{cases} \dfrac{P^{\lambda}-1}{\lambda}, & \lambda\neq 0,\\[0.7em] \ln P, & \lambda = 0,\end{cases}
\qquad
\ell_p(\lambda) = -\frac{n}{2}\ln \hat\sigma^2_g(\lambda) + (\lambda-1)\sum_{t}\ln P_t ,
\end{equation}
where $\hat\sigma^2_g(\lambda)$ is the residual variance of the differenced,
transformed series. The maximiser is $\hat\lambda = 0.42$ with a 95\% likelihood
interval $[0.31,\,0.54]$ (Figure~\ref{fig:boxcox}). Two conclusions follow, and
they point in slightly different directions.

\begin{figure}[htbp]
\centering
\includegraphics[width=\textwidth]{\figdir/f03_boxcox.png}
\caption{(a) Box--Cox profile likelihood for the CPI level. (b) The level is
strongly non-normal, so a transformation is required.}
\label{fig:boxcox}
\end{figure}

Statistically, $\lambda=0$ (the logarithm) lies \emph{outside} the 95\% interval,
so a purist would choose something closer to the square root. Practically, the
likelihood is extremely flat: the profile changes by less than one unit between
$\lambda=0$ and $\lambda=0.42$, and every subsequent conclusion of the analysis
--- the order of differencing, the ACF/PACF patterns, the selected orders, the
residual diagnostics and the forecast accuracy --- is unchanged under $\ln(\cdot)$
or $\sqrt{\cdot}$. The logarithm is retained for three reasons: it linearises the
trend (the log index grows almost exactly linearly, whereas the square root does
not), it maps the series onto a growth rate, and a growth rate is what a
decision-maker budgets for. Table~\ref{tab:transforms} documents the choice.

\input{\tabdir/t03_transformations}

\subsection{Stationarity: ACF, PACF and unit-root tests}
\label{sec:acf}

The autocorrelation function of $y_t$ decays with extraordinary slowness: it is
$0.997$ at lag 1 and still $0.965$ at lag 12 (Table~\ref{tab:acfpacf} and
Figure~\ref{fig:acf}). The partial autocorrelation, in contrast, is essentially
zero after the first lag ($-0.005$ at lag 12). This combination is the textbook
signature of a non-stationary series: the ACF of an AR(1) with $\phi=0.997$ would
look identical, and the PACF tells us that a single AR term would already
``explain'' the level, which is exactly what a unit root does.

\input{\tabdir/t04_acf_pacf}

Formal confirmation uses the two tests of Chapter~2 in parallel, because their
nulls are opposite and the agreement between them is more informative than either
test alone. The augmented Dickey--Fuller statistic for $y_t$ is $-1.29$
($p=0.64$) with a constant and $-1.07$ ($p=0.93$) with a constant and trend, so
the unit-root null cannot be rejected. The KPSS statistic is $4.26$ against a 5\%
critical value of $0.46$, so stationarity is decisively rejected. Both tests agree
that $y_t$ is $I(1)$.

After one difference, $\Delta y_t=\pi_t$, the ACF drops to $0.555$ at lag 1,
$0.383$ at lag 2 and $0.264$ at lag 3, and the ADF statistic ranges from $-8.46$
($p<0.001$) with three augmentation lags to $-2.91$ ($p=0.044$) with twelve, so the
unit root is rejected at conventional lag lengths, while the KPSS test on the same
series is borderline. The interpretation of that borderline result is important and
is taken up in Section~\ref{sec:diag}: the full-sample KPSS statistic on $\pi_t$
is contaminated by the 1968--1982 regime shift in the \emph{mean} of inflation, not
by a stochastic trend. Two facts support this reading. First, over 1980--2020 ---
the post-Volcker sample --- the ADF statistic on $\pi_t$ is $-4.94$ ($p<0.001$),
against $-0.62$ ($p=0.87$) over 1957--1979. Second, the residuals of the finally
selected model pass the Ljung--Box test at every lag up to 24, which could not
happen if $\pi_t$ retained a unit root.

\input{\tabdir/t02_unit_root}

The order of differencing is therefore $d=1$. For the seasonal order, the ACF of
$\pi_t$ shows significant spikes at lags 11, 12 and 13 ($0.42$, $0.45$ and $0.34$,
all far outside the $\pm0.072$ Bartlett band), and the Ljung--Box statistic at lag
12 is $1125$ on 12 degrees of freedom. That is a strong annual echo. It is
\emph{not}, however, evidence that the series needs seasonal differencing: the ADF
test on the annually differenced series gives $-3.01$ ($p=0.034$), no better than
the monthly difference, and seasonal differencing would over-difference a series
whose annual pattern is already partly removed by the source's seasonal adjustment.
The modelling decision is accordingly $D=0$ with a seasonal ARMA block to absorb
the echo --- the specification choice examined in Section~\ref{sec:id}.

\begin{figure}[htbp]
\centering
\includegraphics[width=\textwidth]{\figdir/f04_acf_pacf.png}
\caption{ACF and PACF used for identification. The level (top) has an ACF that
decays almost linearly, the signature of $d=1$. The monthly difference (middle)
has an ACF that dies out within a few lags but retains a peak at lag 12, which is
the seasonal signature. The annual difference (bottom) is shown as a check.}
\label{fig:acf}
\end{figure}

Finally, Figure~\ref{fig:varstab} reports the rolling standard deviation of the
level, of the log level and of the monthly difference over 60-month windows. The
level's variance more than triples across the sample; the log's variance is far
more stable; the differenced series has an essentially constant variance apart from
the 1973--1982 episode. This is the graphical counterpart of the Box--Cox result and
of the stationarity tests.

\begin{figure}[htbp]
\centering
\includegraphics[width=\textwidth]{\figdir/f05_variance_stability.png}
\caption{Rolling 60-month standard deviation. A variance-stabilising transform is
required for the level; the monthly difference is heteroskedastic only around the
two oil shocks.}
\label{fig:varstab}
\end{figure}

% ===========================================================================
\section{Identification and estimation (Chapters 3, 4, 6 and 7)}
\label{sec:id}

\subsection{The multiplicative seasonal model}

Chapter~7 generalises the ARIMA model to the multiplicative seasonal form
\begin{equation}
\underbrace{\phi(B)\Phi(B^{s})}_{\text{AR}}\;
\underbrace{(1-B)^{d}(1-B^{s})^{D}}_{\text{differencing}}\;
y_t \;=\; c \;+\;
\underbrace{\theta(B)\Theta(B^{s})}_{\text{MA}}\;\eps_t ,
\qquad \eps_t \sim WN(0,\sigma^2),
\label{eq:sarima}
\end{equation}
with $B$ the backshift operator, $s=12$ for monthly data, and
$\phi(B)=1-\phi_1B-\dots-\phi_pB^p$, $\theta(B)=1+\theta_1B+\dots+\theta_qB^q$,
and the seasonal polynomials defined analogously in $B^{12}$. With $d=1$ and
$D=0$ established in Section~\ref{sec:acf}, the identifying task is to choose
$p,q$ and the seasonal orders $P,Q$.

\subsection{Identification from the ACF and PACF}

Chapter~6 prescribes reading the orders from the ACF/PACF of the stationary series.
For $\pi_t$ the ACF decays geometrically ($0.555$, $0.383$, $0.264$, $0.15$,
$0.12$, $0.08$, \dots) with a superimposed peak at lag 12, and the PACF has a large
first spike ($0.555$) and a smaller one at lag 12 ($0.161$). A geometric ACF with a
single dominant PACF spike is the signature of an AR(1) or, when the first MA
coefficient is also substantial, an ARMA(1,1); a single spike at the seasonal lag
in the PACF suggests $P=1$. The tentative model is therefore
ARIMA$(1,1,1)(1,0,1)_{12}$, with low-order alternatives
ARIMA$(1,1,1)$, ARIMA$(2,1,1)$ and ARIMA$(1,1,2)$ to be checked.

Identification from the correlogram is a starting point, not an answer. The
theoretical ACF of a model with a seasonal MA term has an infinite tail, and
sample autocorrelations of a series with 767 observations are estimated with
substantial noise. The formal choice is therefore made over a grid, with the
information criteria as the criterion and the residual diagnostics as a filter ---
which is the procedure Chapter~6 actually recommends.

\subsection{Candidate grid and selection}

Every specification in the grid
\begin{equation}
(p,q)\in\{0,1,2,3\}^2,\qquad (P,Q)\in\{(0,0),(1,0),(0,1),(1,1)\},\qquad d=1,\ D=0,
\end{equation}
was estimated by Gaussian maximum likelihood using the Kalman filter, subject to
$k\le 10$ estimated parameters, and each was screened with the Ljung--Box test at
12 and 24 lags. Table~\ref{tab:selection} reports the best of each family;
Figure~\ref{fig:selection} ranks the fourteen best overall.

\input{\tabdir/t05_model_selection}

Three results stand out, and each is a lesson rather than a technicality.

First, a purely \emph{non-seasonal} model is not adequate. The best of that family,
ARIMA$(1,1,3)$ with $d=1$ and no seasonal block, has a BIC of $264.7$ --- a full
$100$ points worse than the selected model --- and its residual autocorrelation at
lag 12 is $+0.213$, far outside the $\pm0.071$ band. Its Ljung--Box $p$-value at
12 lags is below $0.001$: the annual echo survives, and the residuals are not white
noise. This is exactly the failure that Chapter~7 warns about when seasonality is
ignored.

Second, adding seasonal differencing is the wrong remedy. When $D=1$ is imposed
on the same non-seasonal orders, ARIMA$(1,1,1)(1,1,1)_{12}$ has a BIC of $201.6$
--- still $38$ points worse than the selected model --- and its residuals fail the
portmanteau test badly: the Box--Pierce statistic at 24 lags is $60.0$ on 24
degrees of freedom ($p=0.0001$), even though its \emph{contemporaneous} residual
autocorrelation at lag 12 happens to be near zero ($+0.001$). Seasonal differencing
applied to a series without a seasonal unit root \emph{creates} a seasonal MA
structure with a coefficient near $-1$, the classic over-differencing signature.
Section~\ref{sec:acf} already reached the same conclusion from the unit-root tests,
so $D=0$ is retained. Figure~\ref{fig:seasonalblock} shows the like-for-like
comparison.

Third, the seasonal block must contain \emph{both} a seasonal AR and a seasonal MA
term. The selected model, ARIMA$(1,1,3)(1,0,1)_{12}$, has AICc $=126.9$ and BIC
$=163.9$. Dropping the seasonal MA term, giving ARIMA$(1,1,3)(0,0,1)_{12}$, raises
BIC to $241.2$ and its 24-lag Box--Pierce $p$-value falls to $0.011$; dropping the
seasonal AR term, giving ARIMA$(1,1,3)(1,0,0)_{12}$, raises BIC to $233.9$. Figure
~\ref{fig:seasonalblock} isolates the mechanism for the extreme case: with the whole
seasonal block removed, the residual autocorrelation at lag 12 is $+0.213$, against
$-0.018$ for the selected model.

The selection is therefore not "lowest BIC" in a mechanical sense. BIC alone would
prefer ARIMA$(2,1,1)(1,0,1)_{12}$, whose BIC of $162.3$ is $1.6$ points lower;
that model is rejected because its Ljung--Box $p$-value at 12 lags is $0.026$, so
its residuals are not white noise at the 5\% level. Adding the third MA term and
the AR term of the selected model costs $1.6$ BIC points and buys residuals that
pass the portmanteau test at 12 and 24 lags, which is the trade Chapter~6
recommends.

\begin{figure}[htbp]
\centering
\includegraphics[width=\textwidth]{\figdir/f06_model_selection.png}
\caption{Fourteen best specifications by BIC. The vertical axis is truncated so
that the two families (with and without a seasonal block) are both visible; the
selected model is at the top.}
\label{fig:selection}
\end{figure}

\begin{figure}[htbp]
\centering
\includegraphics[width=\textwidth]{\figdir/f07_seasonal_block.png}
\caption{Residual autocorrelation with and without the seasonal MA term. Without
$\Theta(B^{12})$ the annual echo reappears in the residuals at lag 12; with it the
ACF is flat inside the Bartlett bands.}
\label{fig:seasonalblock}
\end{figure}

\subsection{The three estimation methods of Chapter 4}

Chapter~4 requires the method of moments, least squares and maximum likelihood.
All three are applied to the same series, $\pi_t$ (Table~\ref{tab:estimators}),
and the comparison is instructive.

\input{\tabdir/t06_estimators}

\textbf{Method of moments.} For an AR(1), the Yule--Walker equations give
$\hat\phi_1 = r_1$, and for a general AR(p) they give
$\bm\Gamma_p\bm\phi=\bm\gamma_p$. Applied to $\pi_t$: $r_1=0.5558$, so
$\hat\phi_1 = 0.5558$ (Burg's recursive algorithm gives $0.5552$, and the
Yule--Walker solution $0.5551$ --- the three agree to two decimals). The innovation
variance follows from $\hat\sigma^2 = \hat\gamma_0(1-\hat\phi_1^2)$.

For an MA(1) the moments give
$\rho_1 = -\theta/(1+\theta^2)$, so $\theta$ solves the quadratic
$\theta^2+\theta/r_1+1=0$. Here $r_1=0.5558$, and the discriminant
$1/r_1^2-4 = -0.763$ is \emph{negative}: \textbf{the method of moments has no real
solution.} This is not a numerical accident. A moment estimator based on the
first autocorrelation is feasible only for $|r_1|\le 0.5$, and this series has
persistence well above that threshold. The failure is a genuine limitation of the
classical method, and it is the reason the course moves on to least squares and
maximum likelihood. The applied lesson is that moment estimators are convenient
for pure AR models and inadequate for MA components with high persistence.

\textbf{Least squares.} Conditional least squares minimises
$S(\bm\beta)=\sum\hat\eps_t^2$ with $\hat\eps_t = x_t-\mu-\theta\hat\eps_{t-1}$
and $\hat\eps_0=0$, which is nonlinear in $\theta$. A quasi-Newton search gives
$\hat\theta = 0.411$ and $\hat\sigma^2_{\text{CSS}} = 0.0965$ for the MA(1), and
$\hat\phi=0.719$, $\hat\theta=-0.244$ for the ARMA(1,1). Hannan--Rissanen --- the
two-step estimator that regresses $x_t$ on its own lags to obtain residuals and
then uses those residuals as regressors, i.e.\ least squares on an
\emph{estimated} MA term --- gives $\hat\phi = 0.739$ and $\hat\theta=-0.276$.

\textbf{Maximum likelihood.} The exact Gaussian likelihood, evaluated through the
Kalman filter, gives $\hat\theta = 0.4108$ and $\hat\sigma^2 = 0.0964$ for the
MA(1) --- indistinguishable from conditional least squares --- and
$\hat\phi=0.7178$, $\hat\theta=-0.2432$ for the ARMA(1,1), within $0.03$ of the
Hannan--Rissanen values and $0.002$ of conditional LS. The agreement is the
expected outcome for $n=767$: conditional and exact likelihood differ only in the
treatment of the first few observations, so their relative difference is of order
$1/n$. The practical recommendation that follows from Table~\ref{tab:estimators}
is that for samples of this size the three methods are interchangeable, and the
choice should be made on convenience; ML is used for the seasonal model because it
handles the seasonal polynomials and the differencing in one consistent
likelihood.

\subsection{The selected model}

Table~\ref{tab:params} reports the ML estimates. The fitted model is
\begin{equation}
(1-0.9588 B)(1-0.9747 B^{12})\,\Delta y_t
= 0.0003 + (1-0.5823B-0.1058B^2-0.1281B^3)(1-0.8676 B^{12})\,\hat\eps_t ,
\label{eq:fitted}
\end{equation}
with $\hat\sigma_\eps = 0.2586$, log-likelihood $-55.36$, AIC $126.71$, AICc
$126.90$ and BIC $163.85$.

\input{\tabdir/t07_parameters}

Four features of these estimates deserve comment.

\emph{Stationarity and invertibility.} The inverse AR roots have moduli $0.754$ and
$0.244$ (the second from $\phi_1$ and $\Phi_1$ combined), all below one, so the
model is stationary; the inverse MA roots have moduli $0.62$ (from $\theta_1$) and
$0.87$ (from $\Theta_1$), all below one, so it is invertible. The seasonal AR
coefficient $\hat\Phi_1 = 0.9747$ is close to the unit circle: the annual component
is highly persistent, with a seasonal half-life of about $12\ln 0.5/\ln 0.975
\approx 330$ months, which is the formal statement of what Figure~\ref{fig:stl}
showed visually --- this is a slow-moving annual pattern, not a calendar effect
that resets each year.

\emph{Persistence.} The largest inverse AR modulus is $0.754$, so the half-life of
a shock to the log price level is $\ln 0.5/\ln 0.754 = 16.5$ months. A one
percentage-point surprise in monthly inflation therefore leaves about half a
percentage point in the annual rate after eighteen months. Inflation is persistent
but mean-reverting on a horizon that matters for a one- to three-year contract.

\emph{The constant.} The estimated constant is $\hat c = 0.0003$ with a standard
error of $0.0003$ and a $p$-value of $0.39$, so it is not statistically
significant. Dropping it is nevertheless a mistake, and Section~\ref{sec:fcst}
shows why in out-of-sample terms: a constant of exactly zero in
The seasonal model above forces the long-run drift of the log price level, and
hence long-run inflation, to be exactly zero, which contradicts both the data and
economics. Retained, the implied long-run inflation rate is
$\hat c/[(1-\sum\hat\phi)(1-\hat\Phi_1)] = 0.0003/(0.0412\times0.0253)$, i.e.\
$0.243\%$ per month or about $2.9\%$ a year --- within a whisker of the Federal
Reserve's 2\% target plus the historical bias of the CPI.

\emph{The MA(3) term.} Three MA coefficients, not one, are required, and all three
are significant at the 1\% level. The need for $\theta_3$ reflects the short-run
dynamics of monthly inflation: price changes are not a first-order process but
reverberate for about a quarter, consistent with the staggered contract structure
of the US economy.

Figure~\ref{fig:fitted} shows that the fitted values track the observed log index
closely, and that the residuals are centred on zero with a fairly constant spread,
apart from the two known episodes.

\begin{figure}[htbp]
\centering
\includegraphics[width=\textwidth]{\figdir/f08_fitted_residuals.png}
\caption{Observed and fitted log CPI (top) and residuals (bottom). The model
reproduces the trend and the cyclical swings; the largest residuals are the 1973
oil embargo and the 2008 financial crisis.}
\label{fig:fitted}
\end{figure}

% ===========================================================================
\section{Diagnostic checking (Chapter 6, stage 3)}
\label{sec:diag}

Chapter~6 lists four residual checks --- zero mean, constant variance, absence of
residual autocorrelation, and adequacy relative to neighbouring models. Each is
reported in Table~\ref{tab:diagnostics} and discussed in turn. The model is
re-estimated on the full sample for this section ($n=768$, $n=767$ usable
residuals after the diffuse initialisation is dropped).

\input{\tabdir/t12_residual_diagnostics}

\subsection{Zero mean}

The residual mean is $0.00018$ percentage points with a standard error of
$0.00939$, so $t=0.02$ and $p=0.985$: the null of a zero mean is not remotely
rejected, and the 95\% interval $[-0.018,\,0.019]$ is tight relative to
$\hat\sigma_\eps=0.259$. There is no systematic bias left in the conditional mean.

\subsection{Constant variance}

The residuals are split into two halves of 383 observations. The variances are
$0.06779$ and $0.06747$, a ratio of $1.005$ with an $F$-test $p$-value of $0.963$,
so the unconditional variance is constant across the sample. Regressing the squared
residuals on a time trend gives a slope of $2\times10^{-5}$ with $t=0.55$
($p=0.58$), confirming that the variance does not drift with time.

This is where the honest qualification belongs. The \emph{unconditional} variance
is constant, but the \emph{conditional} variance is not: the ARCH-LM statistic at
12 lags is $49.6$ ($p<0.001$) and the Ljung--Box test on the squared residuals
gives $76.4$ ($p<0.001$). Volatility clusters --- the early 1970s, the Volcker
years and 2008--2009 are visibly more volatile than the mid-1990s --- and a model
with homoskedastic innovations cannot represent that. The consequence for the
forecast is precise and worth stating carefully: the point forecasts are
unaffected, because they depend only on the conditional mean equation, but the
\emph{Gaussian confidence intervals are too narrow in turbulent periods and
slightly too wide in calm ones}. A GARCH extension (Tsay, 2010, ch.~3) would fix
the intervals; it is outside the scope of this course's Box--Jenkins cycle, and
the intervals reported in Section~\ref{sec:fcst} are correctly labelled as
conditional on homoskedasticity.

\subsection{Residual autocorrelation}

This is the decisive check. The residual ACF and PACF (Figure~\ref{fig:resacf})
mostly stay within the $\pm0.071$ Bartlett band; the largest absolute
autocorrelation over the first twelve lags is $0.098$, at lag 11, and the next
largest is well inside the band. The Ljung--Box statistics
are $3.27$ at 6 lags ($p=0.77$), $15.66$ at 12 ($p=0.21$), $28.36$ at 18 ($p=0.06$)
and $33.10$ at 24 ($p=0.10$), and the Box--Pierce statistic at 24 lags is $32.44$
($p=0.12$). None of these rejects white noise.

One honest caveat: at 36 lags the Ljung--Box $p$-value falls to $0.042$. That
marginal rejection must be read in context. A portmanteau statistic computed at 36
lags accumulates 36 sample autocorrelations; with $n=767$ and no individual
autocorrelation outside the band, the test is picking up the small accumulated
noise rather than a specific structure, and it is being applied at a lag length
far beyond the estimated orders (the rule of thumb is to test at
$m\approx\ln n$ to $2\ln n$, i.e.\ 7 to 13). It is reported rather than hidden, but
it does not outweigh the evidence at 12 and 24 lags.

\begin{figure}[htbp]
\centering
\includegraphics[width=\textwidth]{\figdir/f15_residual_acf_pacf.png}
\caption{Residual ACF and PACF with Bartlett bands. The shaded region marks lags
11--13, where the seasonal block acts; there is no spike there.}
\label{fig:resacf}
\end{figure}

\subsection{Adequacy and parsimony}

Table~\ref{tab:neighbours} compares the selected model with its immediate
neighbours, all estimated on the full sample. Removing the AR term is catastrophic
(AICc rises from $126.9$ to $167.8$ and the residuals fail Ljung--Box with
$p<0.001$). Removing the seasonal AR term is worse (AICc $208.8$). Adding a fourth
MA term improves the AICc by only $0.4$ and the BIC not at all, so the extra
parameter is not justified. The selected model is therefore the most parsimonious
adequate specification in its neighbourhood, which is exactly the criterion of
Chapter~6.

\input{\tabdir/t13_neighbouring_models}

\subsection{Influence of the largest residuals}

The six largest standardised residuals are 2008-11 ($z=-6.13$), 1973-08
($z=6.10$), 2008-10 ($z=-3.85$) and three months of the 1980 credit-control
episode. Re-estimating the model with the two largest observations deleted changes
the coefficients by at most $0.058$ --- $22\%$ of one innovation standard deviation
--- and leaves every parameter's sign and significance unchanged. The estimates are
not driven by outliers; the outliers are genuinely unusual events that the model
correctly refuses to fit.

\subsection{Residual normality}

The Jarque--Bera test rejects normality overwhelmingly ($p=6.9\times10^{-119}$),
with skewness $-0.25$ and excess kurtosis $4.13$; the normal probability plot
(Figure~\ref{fig:qq}) shows the characteristic S-shape of a fat-tailed
distribution. This matters for the forecast intervals: with excess kurtosis, the
$t$-statistic $\pm1.96$ does not correspond exactly to a 95\% probability, and the
true coverage of a nominal 95\% interval built this way is slightly below 95\%.
The interval is still reported in the standard way --- this is the Gaussian
likelihood's own interval --- but the qualification belongs in the interpretation,
and Section~\ref{sec:interpretation} reports the realised coverage as a check.

\begin{figure}[htbp]
\centering
\includegraphics[width=\textwidth]{\figdir/f16_qq_and_squared.png}
\caption{(a) Normal probability plot of the residuals: the tails are much heavier
than Gaussian. (b) Autocorrelation of the squared residuals: volatility clusters.}
\label{fig:qq}
\end{figure}

\begin{figure}[htbp]
\centering
\includegraphics[width=\textwidth]{\figdir/f14_residuals.png}
\caption{Standardised residuals over time and their distribution. Three
observations exceed three standard deviations; all three are dated economic
events.}
\label{fig:resid}
\end{figure}

% ===========================================================================
\section{Forecasting and out-of-sample evaluation (Chapters 5--6)}
\label{sec:fcst}

\subsection{The MMSE forecast and its error variance}

Chapter~5 derives the minimum-mean-square-error forecast by conditional
expectation. Write the model in its MA$(\infty)$ form, $y_t = \sum_{j\ge0}\psi_j
\eps_{t-j}$, with $\psi_0=1$. Then the forecast is the conditional expectation,
\begin{equation}
\yhat_{t+h|t} \;=\; \E[y_{t+h}\mid y_t,y_{t-1},\dots] \;=\; \sum_{j=h}^{\infty}\psi_j\,\eps_{t+h-j},
\qquad
e_{t+h|t} = \sum_{j=0}^{h-1}\psi_j\,\eps_{t+h-j},
\end{equation}
so it is obtained by setting future innovations to their expectation of zero.
The forecast error variance is
\begin{equation}
\Var\!\left(e_{t+h|t}\right) \;=\; \sigma^2\sum_{j=0}^{h-1}\psi_j^{2},
\qquad
\yhat_{t+h|t}\;\pm\;z_{\alpha/2}\,\hat\sigma_\eps\sqrt{\textstyle\sum_{j=0}^{h-1}\hat\psi_j^2}
\label{eq:interval}
\end{equation}
the second line of which is the Gaussian confidence interval. The $\psi$-weights are the impulse responses
of the fitted \emph{integrated} model, which for the model of
Table~\ref{tab:params} begin
$\psi_0=1$, $\psi_1=1.359$, $\psi_2=1.619$, $\psi_3=1.740$, $\psi_4=1.856$, and by
$h=12$ have flattened to $\psi_{11}=2.550$.

The recursion was implemented independently and checked against the Kalman filter's
own interval: the two agree to $6\times10^{-8}$ at every horizon (Table~\ref{tab:forecastse}),
which validates both the $\psi$-recursion and the fit. The important economic
point is in the second row of that table. Because the ARMA terms are
mean-reverting, the model's forecast error at $h=60$ is $\pm8.62$, whereas a
random walk's would be $\sigma\sqrt{60}=\pm2.00$ in \emph{per-month} units, i.e.\
the model's interval is about three times tighter after five years than the naive
benchmark's. Figure~\ref{fig:fvar} shows how the gap opens with the horizon.

\begin{figure}[htbp]
\centering
\includegraphics[width=\textwidth]{\figdir/f10_forecast_variance.png}
\caption{Forecast standard error by horizon: the $\psi$-weight recursion against
the random-walk benchmark $\sigma\sqrt{h}$.}
\label{fig:fvar}
\end{figure}

\input{\tabdir/t11_forecast_se}

\subsection{Hold-out design}

The evaluation holds out the last five years, 2016-01 to 2020-12 ($h=60$), and
re-estimates the model on 1957-01 to 2015-12 ($n=708$). This is a genuine
out-of-sample test: no information after 2015-12 enters the point forecasts. Three
benchmarks are used, each of which a practitioner might use in the absence of a
model:

\begin{itemize}[leftmargin=1.4em,itemsep=0.1em]
\item \textbf{Random walk}: $\yhat_{T+h|T}=y_T$, the forecast implied by
``inflation does not change''.
\item \textbf{Random walk with drift}: $\yhat_{T+h|T}=y_T+h\hat\mu$, with $\hat\mu$
the mean log change of the training sample.
\item \textbf{Seasonal random walk}: $\yhat_{T+h|T}=y_{T+h-12}$, the forecast
implied by ``the same month next year''.
\end{itemize}

\subsection{Results}

Table~\ref{tab:accuracy} reports the four standard criteria on the log index, and
Figure~\ref{fig:fan} shows the forecast against the realised path.

\input{\tabdir/t08_accuracy}

\begin{figure}[htbp]
\centering
\includegraphics[width=\textwidth]{\figdir/f09_forecast_fanchart.png}
\caption{Out-of-sample forecast with 80\% and 95\% intervals. The realised path
stays inside the 95\% band throughout; the random-walk benchmark (dotted) sits well
below it.}
\label{fig:fan}
\end{figure}

The model dominates every benchmark. Its RMSE on the log index is $1.79$ against
$6.23$ for the random walk, $4.49$ for the random walk with drift and $1.86$ for
the seasonal random walk --- a $71\%$ reduction relative to the naive benchmark and
a small but consistent improvement on the seasonal one. The MASE of $0.011$ is
consistent with a model that beats the seasonal random walk out of sample. On the
index itself (Table~\ref{tab:accuracylevel}) the MAPE over five years is $1.59\%$,
compared with $5.37\%$ for the random walk.

\input{\tabdir/t09_accuracy_level}

The comparison between the two tables carries a methodological warning worth
stating. On the index, \emph{every} model has a MAPE below 6\%, including the
random walk, because the level is dominated by the trend. Judging a model by its
level forecast is therefore almost guaranteed to produce flattering numbers.
Accuracy must be judged on the stationary transformation, which is why
Table~\ref{tab:accuracy} is the primary evidence and Table~\ref{tab:accuracylevel}
is reported only for interpretability.

Figure~\ref{fig:fan} reveals a systematic feature of the forecast, not an error:
the model's path runs slightly \emph{below} the realised one after 2018. The
realised 2020 inflation of $1.35\%$ is below the model's long-run $2.9\%$, but the
log \emph{index} in 2020 ended higher than predicted because the model's estimated
cyclical position implied reversion towards the mean that did not occur. The
consequence is that the 60-month-ahead point forecast is biased upward by about
$1.4$ log points, while the interval is wide enough to contain the truth --- which
is the correct behaviour of an honest forecast.

\subsection{Recursive updating and the rolling evaluation}

Chapter~5 also requires forecasts to be updated as new data arrive. Figure~\ref{fig:update}
reports the same 36-month-ahead forecast made at four successive origins; the
paths are revised as information accrues, and the revisions are largest in the
first few months after each origin, exactly as the $\psi$-weights imply.

\begin{figure}[htbp]
\centering
\includegraphics[width=\textwidth]{\figdir/f12_forecast_updating.png}
\caption{Recursive updating: the same multi-step forecast made at four origins.}
\label{fig:update}
\end{figure}

A one-step-ahead rolling evaluation over the entire hold-out --- re-estimating the
model at each of the 60 origins and forecasting one month ahead --- gives an RMSE
of $0.225$, an MAE of $0.173$ and a bias of $+0.007$. The theoretical innovation
standard deviation is $\hat\sigma_\eps=0.261$, so the realised one-step error is
\emph{smaller} than the theoretical value and the bias is negligible. This is the
strongest evidence in the report that the conditional mean is correctly specified:
if a term were missing, the one-step errors would be systematically larger than
$\sigma$.

\subsection{Sensitivity checks}

Two sensitivity checks are reported in full honesty.

\textbf{Dropping the constant.} The information criteria marginally prefer the
model without the constant (AICc $125.7$ against $126.9$), which is the expected
outcome for a parameter with a $p$-value of $0.39$. On the hold-out, however, the
constant earns its place decisively: RMSE $1.79$ with it, $3.40$ without, and the
paths diverge by $3.2$ log points at $h=60$. The economic reason is that a zero
drift imposes zero long-run inflation, so the model without the constant forecasts
a flat price level over five years --- a restriction the data reject decisively.
Retaining an economically motivated parameter that is statistically insignificant
is a legitimate identification choice in Box--Jenkins practice, and this is the
quantitative justification.

\textbf{Annual inflation forecasts.} For a decision-maker the most useful output
is not the log index but the annual inflation rate implied by the forecast path,
$\hat\pi^{12}_{t+h|t}=\hat y_{t+h|t}-y_{t+h-12}$ (observed history for $h\le12$,
model-based otherwise). Over the hold-out these conditional forecasts have an RMSE
of $0.76$ percentage points against realised annual inflation, with a mean forecast
of $1.38\%$ against a mean realisation of $1.76\%$ (Figure~\ref{fig:anninfl}). The
forecast undershoots the 2017 recovery slightly and tracks the 2019--2020 slowdown
well.

\begin{figure}[htbp]
\centering
\includegraphics[width=\textwidth]{\figdir/f13_annual_inflation_forecast.png}
\caption{Annual inflation implied by the forecast path against the realised
twelve-month rate.}
\label{fig:anninfl}
\end{figure}

\begin{figure}[htbp]
\centering
\includegraphics[width=\textwidth]{\figdir/f11_accuracy_comparison.png}
\caption{RMSE and MAE on the log index, model against the three naive benchmarks.}
\label{fig:accuracy}
\end{figure}

% ===========================================================================
\section{Interpretation in decision-relevant language}
\label{sec:interpretation}

\subsection{What the model says}

Translating the estimates into the terms in which the original question was posed:

\begin{enumerate}[leftmargin=1.6em,itemsep=0.35em]
\item \textbf{The trend rate.} US consumer price inflation has averaged about
$2.9\%$ a year over 1957--2020 once the short-run noise is filtered out. That is
the number to put in a long-horizon contract, and it is remarkably close to the
Federal Reserve's 2\% target plus the known upward bias of the CPI relative to
the Fed's preferred PCE measure.

\item \textbf{How fast surprises fade.} A surprise in monthly inflation has a
half-life of about 16 months. For a treasury desk this is the operational number:
a 1 percentage-point surprise persists in the annual rate for roughly a year and a
half, so it should be treated as a two-year problem rather than a one-month one.

\item \textbf{The annual pattern.} The seasonal AR coefficient of $0.975$ implies a
seasonal half-life of about 27 years --- the annual component is essentially a
slow-moving feature of the inflation process, not a calendar effect that resets.
Practically, this means that the month of the year is almost irrelevant to a
forecast, but the \emph{phase} of the multi-year inflation cycle is not.

\item \textbf{The uncertainty.} The 95\% interval around the 12-month-ahead
forecast is $\pm1.81$ log points, i.e.\ roughly $\pm1.8$ percentage points on the
annual rate; at 60 months it is $\pm8.62$ log points. For comparison, the
random-walk interval at 60 months is about three times wider. This is the concrete
value of the model: it converts a five-year inflation exposure from a
near-unknowable quantity into one with a usable range.
\end{enumerate}

\subsection{What it means for a business decision}

Three concrete uses follow.

\textbf{Pricing a dollar-denominated loan.} A Tunisian bank lending in dollars for
five years needs a nominal compensation for expected US inflation plus a risk
premium for inflation uncertainty. The model supplies both numbers: $2.9\%$ for the
first and, from the interval formula above, a 95\% band of roughly
$\pm1.8$ points at one year widening to $\pm8.6$ at five. A prudent contract prices
at the upper part of that band rather than at the mean.

\textbf{Indexation clauses.} A supply contract indexed to the CPI needs a forecast
and a tolerance. The annual-rate forecast has an RMSE of $0.76$ percentage points
over the 60 hold-out months; a tolerance of $\pm1.5$ points would have covered the
realised path in 58 of those 60 months, and $\pm1.0$ point in 47.

\textbf{Scenario planning.} Because the model is explicitly a model of the
\emph{conditional mean}, its residuals tell the risk manager where the risk lies.
The two largest residuals (1973-08 and 2008-11) are supply shocks and financial
shocks respectively --- precisely the events that a univariate extrapolation cannot
anticipate. A scenario exercise should therefore stress the model's central path
with an oil-price shock of the 1973 magnitude and a demand collapse of the 2008
magnitude, rather than perturbing the ARMA parameters.

\subsection{The comparison a manager should actually make}

The RMSE comparison in Table~\ref{tab:accuracy} is worth restating in money terms.
Over a five-year horizon the random walk is wrong by $6.23$ log points on average,
which is about $6.4\%$ of the price level; the model is wrong by $1.79$, or about
$1.8\%$. On a \$10 million five-year dollar exposure, that is the difference between
a typical pricing error of \$640,000 and one of \$180,000. The model does not
eliminate the exposure, but it reduces the expected pricing error by roughly
three-quarters.

% ===========================================================================
\section{Limitations}
\label{sec:limits}

Five limitations should be stated plainly, because each of them bounds how the
results may be used.

\textbf{No causal content.} A univariate ARIMA model extrapolates the process that
generated the past. It cannot anticipate a change in monetary policy, a supply
shock or a pandemic. This is not a defect of this particular specification but a
property of the class of models, and the 2020 observations in the sample
(monthly inflation falling to $0.002$ points, essentially zero, in May 2020) are a reminder that it bites.

\textbf{Regime instability.} The sample contains at least three regimes: the
post-war dollar-gold system with inflation near $2\%$, the Great Inflation with
inflation near $7\%$, and the Great Moderation with inflation near $2.5\%$. A
single set of parameters is a compromise across these regimes. The large residual
at 1973-08 and the marginal Ljung--Box rejection at 36 lags are the visible
symptoms. A Markov-switching or threshold model would be the natural extension.

\textbf{Heteroskedastic residuals.} The ARCH-LM test rejects conditional
homoskedasticity, so the Gaussian confidence intervals are too narrow in turbulent
periods. The point forecasts are unaffected, but a risk application should use a
GARCH-augmented interval.

\textbf{The log transform is a compromise.} The Box--Cox estimate of $0.42$ lies
outside the 95\% interval for $\lambda=0$. The analysis was checked under the
square-root transform and nothing material changes, but the choice is a
convenience rather than a likelihood result.

\textbf{An insignificant constant is retained.} This is deliberate and defended in
Section~\ref{sec:fcst} by out-of-sample performance, but a reader applying the
strict information-criterion rule of Chapter~6 would drop it, and should know that
the resulting model forecasts a flat price level over five years.

% ===========================================================================
\section{Conclusion}
\label{sec:conclusion}

The project set out to answer a practical question --- what US inflation rate
should be budgeted for, and how wide is the interval around it --- using the full
Box--Jenkins methodology, and it answered it with a specific, validated model.

Following the course's own sequence: Chapter~1's components were identified from
the plot (a dominant log-linear trend, long cycles, a weak seasonal echo, and
irregular shocks in 1973 and 2008); Chapter~2's transformations and tests fixed
$d=1$, $D=0$ with the Box--Cox profile, the ACF/PACF decay and the agreement of the
ADF and KPSS tests; Chapter~4's three estimation methods were compared and found to
agree to two decimals, with the method of moments failing outright for the MA(1)
because the first autocorrelation exceeds the feasibility threshold; Chapters~6
and~7's multiplicative seasonal specification was selected over 51 alternatives,
with both the non-seasonal and the seasonally differenced families rejected on
diagnostic grounds; and Chapter~5's MMSE recursion produced forecasts whose
$\psi$-weight-based intervals reproduce the Kalman filter's to eight decimal
places.

The result is ARIMA$(1,1,3)(1,0,1)_{12}$ on the log CPI, which passes the
portmanteau, zero-mean and constant-variance checks, beats the random walk by
$71\%$ in RMSE over a five-year hold-out, and implies a long-run inflation rate of
$2.9\%$ a year with a shock half-life of 16 months. Its limitations --- fat tails,
volatility clustering, regime dependence and no causal content --- are the agenda
for any practitioner who wants to go further than a univariate extrapolation.

% ===========================================================================
\clearpage
\appendix
\section{Reproducibility and code}
\label{app:code}

Every number, table and figure in this report is produced by a Python script in the
project repository, in the order below. The scripts write their results to
\texttt{results/*.json} (every reported statistic), \texttt{reports/tables/*.tex}
(the \LaTeX{} tables included in this document) and
\texttt{reports/figures/*.png} (the figures).

\begin{table}[htbp]
\centering
\small
\caption{Pipeline. Each stage is a standalone script; running them in order
regenerates the entire report from the raw Kaggle CSV.}
\label{tab:pipeline}
\begin{tabular}{@{}lll@{}}
\toprule
Script & Purpose & Course chapter \\
\midrule
\texttt{01\_acquire\_data.py} & download and verify the Kaggle file (SHA-256), audit it & --- \\
\texttt{02\_prepare\_data.py} & clean, parse dates, build $P_t,y_t,\pi_t,\pi^{12}_t$ & 1--2 \\
\texttt{03\_components\_stationarity.py} & components, Box--Cox, ACF/PACF, ADF/KPSS & 1--2 \\
\texttt{04\_identify\_estimate.py} & grid identification, MoM/CSS/HR/MLE, final SARIMA & 3, 4, 6, 7 \\
\texttt{05\_forecasting.py} & MMSE forecasts, $\psi$-intervals, hold-out, benchmarks & 5--6 \\
\texttt{06\_diagnostics.py} & zero mean, variance, portmanteau, parsimony & 6 \\
\texttt{07\_build\_report.py} & compile this \LaTeX{} document & --- \\
\bottomrule
\end{tabular}
\end{table}

Reproduction requires Python 3.11 with \texttt{pandas}, \texttt{numpy},
\texttt{scipy}, \texttt{statsmodels} and \texttt{matplotlib}:

\begin{quote}\ttfamily\small
python -m venv .venv \&\& .venv/bin/pip install -r requirements.txt\\
.venv/bin/python run\_all.py
\end{quote}

To obtain the dataset from Kaggle directly rather than from the archived mirror,
set \texttt{KAGGLE\_USERNAME} and \texttt{KAGGLE\_KEY} in the environment (or place
\texttt{kaggle.json} in \texttt{\textasciitilde/.kaggle/}) before running stage 1;
the script then downloads the dataset with \texttt{kagglehub} and verifies it
against the same digest.

\section{Derivations}
\label{app:derivations}

\subsection*{A.1 The Yule--Walker equations and the moment estimator for MA(1)}

For a stationary AR(p), multiplying $x_t=\phi_1x_{t-1}+\dots+\phi_px_{t-p}+\eps_t$
by $x_{t-k}$ and taking expectations gives $\gamma_k=\phi_1\gamma_{k-1}+\dots+
\phi_p\gamma_{k-p}$ for $k\ge1$, i.e.\ $\bm\Gamma_p\bm\phi=\bm\gamma_p$ with
$(\bm\Gamma_p)_{ij}=\gamma_{|i-j|}$. Dividing by $\gamma_0$ gives the sample
estimator $\hat{\bm\phi}$ from the sample autocorrelations.

For an MA(1), $x_t=\mu+\eps_t+\theta\eps_{t-1}$, so $\gamma_0=\sigma^2(1+\theta^2)$,
$\gamma_1=\sigma^2\theta$ and $\gamma_k=0$ for $k\ge2$. Hence
$\rho_1=\theta/(1+\theta^2)$, which inverts to $\theta^2+\theta/\rho_1+1=0$ with
roots $\theta = (-1/\rho_1 \pm \sqrt{1/\rho_1^{2}-4})/2$. Real roots require
$|\rho_1|\le1/2$, and the invertible root is the one with $|\theta|<1$. For
$\pi_t$, $\rho_1=0.5558>1/2$, so the discriminant is negative and the moment
estimator does not exist.

\subsection*{A.2 Forecast error variance from the $\psi$-weights}

With $y_t=\sum_{j\ge0}\psi_j\eps_{t-j}$ and $\eps$ white noise,
$\E[y_{t+h}\mid y_t,\dots]=\sum_{j\ge h}\psi_j\eps_{t+h-j}$, so the forecast error
is $e_{t+h|t}=\sum_{j=0}^{h-1}\psi_j\eps_{t+h-j}$. Since the $\eps$ are
uncorrelated with variance $\sigma^2$,
$\Var(e_{t+h|t})=\sigma^2\sum_{j=0}^{h-1}\psi_j^2$, which is
the interval formula above. For an ARIMA$(p,d,q)$ the $\psi_j$ solve
$\psi_0=1$ and $\psi_k=\theta_k-\sum_{j=1}^{\min(k,p+d)}\phi^{*}_j\psi_{k-j}$, with
$\phi^{*}$ the coefficients of $(1-B)^d\phi(B)$; numerically it is safer to read
them off the fitted model's impulse responses, as the code does.

\subsection*{A.3 Long-run mean of an ARIMA$(p,1,q)$ with a constant}

Because $\Delta y_t$ is stationary with mean $\mu_\Delta = c/\phi(1)\Phi(1)$,
we have $\E[\Delta y_t]=\mu_\Delta$ and, by summing, $\E[y_{t+h}-y_t]=h\mu_\Delta$.
The quantity $\mu_\Delta$ is the long-run \emph{monthly} log inflation rate, so the
annual rate is $12\mu_\Delta$. For the fitted model,
$\phi(1)\Phi(1)=(1-0.9588)(1-0.9747)=0.0412\times0.0253=0.001042$ and
$\mu_\Delta = 0.0003/0.001042 = 0.243\%$ per month, i.e.\ $2.91\%$ per year.

\section*{References}
\begin{thebibliography}{9}
\setlength{\itemsep}{0.15em}

\bibitem[Box et al.(2015)]{box2015}
Box, G.~E.~P., Jenkins, G.~M., Reinsel, G.~C., \& Ljung, G.~M. (2015).
\emph{Time Series Analysis: Forecasting and Control} (5th ed.). Wiley.

\bibitem[Cryer \& Chan(2008)]{cryer2008}
Cryer, J.~D., \& Chan, K.-S. (2008).
\emph{Time Series Analysis: With Applications in R} (2nd ed.). Springer.

\bibitem[Dickey \& Fuller(1979)]{dickey1979}
Dickey, D.~A., \& Fuller, W.~A. (1979). Distribution of the estimators for
autoregressive time series with a unit root. \emph{Journal of the American
Statistical Association}, 74(366), 427--431.

\bibitem[Hannan \& Rissanen(1982)]{hannan1982}
Hannan, E.~J., \& Rissanen, J. (1982). Recursive estimation of mixed
autoregressive-moving average order. \emph{Biometrika}, 69(1), 81--94.

\bibitem[Hyndman \& Athanasopoulos(2021)]{hyndman2021}
Hyndman, R.~J., \& Athanasopoulos, G. (2021).
\emph{Forecasting: Principles and Practice} (3rd ed.). OTexts.

\bibitem[Kwiatkowski et al.(1992)]{kpss1992}
Kwiatkowski, D., Phillips, P.~C.~B., Schmidt, P., \& Shin, Y. (1992). Testing the
null hypothesis of stationarity against the alternative of a unit root.
\emph{Journal of Econometrics}, 54(1--3), 159--178.

\bibitem[Ljung \& Box(1978)]{ljung1978}
Ljung, G.~M., \& Box, G.~E.~P. (1978). On a measure of lack of fit in time series
models. \emph{Biometrika}, 65(2), 297--303.

\bibitem[Osborn et al.(1988)]{ocsb1988}
Osborn, D.~R., Chui, A.~P.~L., Smith, J.~P., \& Birchenhall, C.~R. (1988).
Seasonality and the order of integration for consumption.
\emph{Oxford Bulletin of Economics and Statistics}, 50(4), 361--377.

\bibitem[Said \& Dickey(1984)]{said1984}
Said, S.~E., \& Dickey, D.~A. (1984). Testing for unit roots in autoregressive
moving average models of unknown order. \emph{Biometrika}, 71(3), 599--607.

\bibitem[Tsay(2010)]{tsay2010}
Tsay, R.~S. (2010). \emph{Analysis of Financial Time Series} (3rd ed.). Wiley.

\end{thebibliography}

\end{document}
